\documentclass[11pt]{article}
\usepackage{amsmath,amssymb,amsthm}

\newtheorem{theorem}{Theorem}
\newtheorem{corollary}[theorem]{Corollary}

\newcommand{\Up}{U_+}

\title{Half-coalescence and reset renewal on a record-atom tail}
\author{}
\date{27 August 2026}

\begin{document}
\maketitle

Let
\[
 x_{t+1}=\Up(x_t):=\frac{3x_t+1}{2^{a_t}},
 \qquad
 a_t=v_2(3x_t+1),
 \qquad
 A_t=\sum_{h<t}a_h,
 \tag{1}
\]
be an injective infinite accelerated orbit of positive odd integers.  Put
\[
 Q_t:=\frac{3^t}{2^{A_t}},
 \qquad
 0<\mathcal K:=\lim_{t\to\infty}\frac{x_t}{Q_t}<\infty,
 \tag{2}
\]
and assume that $Q_t\to\infty$.  Let $j_0<j_1<\cdots$ be consecutive
coefficient tail records.  On an interior deficit-four tail, every record
gap is one of the atoms
\[
 A:=1,
 \qquad B:=12,
 \qquad C:=1113,
 \qquad D:=1122.
 \tag{3}
\]
Write
\[
 b_i:=v_3(x_{j_i}+1).
 \tag{4}
\]
The decoder transport at record endpoints is
\[
 A:b\mapsto b+1,
 \qquad B:b\mapsto0,
 \qquad C:b\mapsto1,
 \qquad D:b\mapsto0.
 \tag{5}
\]
The synchronized companion is the sequence determined by
\[
 2^{a_t}e_{t+1}=3e_t-1,
 \qquad e_t>1,
 \qquad \frac{e_t}{Q_t}\longrightarrow0.
\]
Adding its recurrence to (1) shows that
\[
 \frac{x_t+e_t}{Q_t}
\]
is constant; the boundary condition and (2) identify this constant as
$\mathcal K$.  Thus $x_t+e_t=\mathcal KQ_t$ at every state.

\begin{theorem}[exact half-coalescence at the atom $C$]
Suppose that a record transition starts at $x$ and has exponent word
$C=1113$.  Put
\[
 y:=\Up^4(x).
 \tag{6}
\]
Then
\[
 x\equiv79\pmod {128},
 \qquad
 y=\frac{81x+65}{64}.
 \tag{7}
\]
The smaller positive odd integer
\[
\boxed{z:=\frac{x-1}{2}}
\tag{8}
\]
has exact accelerated exponent word $1121$ and coalesces with $x$ at the
same endpoint:
\[
\boxed{\Up^4(z)=y.}
\tag{9}
\]

Let $e_x,e_y>1$ be the synchronized companion values at the two record
endpoints.  Then
\[
 64e_y=81e_x-65.
 \tag{10}
\]
The alternative companion start
\[
 e_z:=\frac{e_x+1}{2}>1
 \tag{11}
\]
is transported to $e_y$ by the same alternative word $1121$, and
\[
\boxed{z+e_z=\frac{x+e_x}{2}.}
\tag{12}
\]

If $x$ belongs to the unbounded orbit of a globally least positive odd
unbounded seed $n$, then
\[
\boxed{x\geq2n+1.}
\tag{13}
\]
If $Q$ is the coefficient at $x$, so that $x+e_x=\mathcal KQ$, then
\[
\boxed{Q>\frac{2(n+1)}{\mathcal K}.}
\tag{14}
\]
\end{theorem}

\begin{proof}
The affine carry of the word $1113$ is $65$ and its total exponent is
$6$, which gives the second identity in (7).  Since $y$ is odd,
\[
 81x+65\equiv64\pmod {128}.
\]
Since $81^{-1}\equiv49\pmod {128}$, it follows that
$x\equiv-49\equiv79\pmod {128}$, proving the first identity in (7).

Starting from $z=(x-1)/2$, the four states generated by the proposed word
$1121$ are
\[
 \frac{x-1}{2},
 \quad
 \frac{3x-1}{4},
 \quad
 \frac{9x+1}{8},
 \quad
 \frac{27x+11}{32},
 \quad
 \frac{81x+65}{64}.
 \tag{15}
\]
The congruence $x\equiv79\pmod {128}$ shows directly that the extracted
powers of two are respectively $1,1,2,1$.  The last state in (15) is
$y$, proving (8)--(9).

Iteration of the negative companion recurrence through $1113$ gives
(10).  The word $1121$ has total exponent $5$ and carry $73$, so its
companion equation is
\[
 32e_y=81e_z-73.
 \tag{16}
\]
Substitution of (11) into the right side of (16), followed by (10), proves
(16).  Equation (12) is immediate from (8) and (11).

The orbit of $z$ reaches $y$ and then follows the unbounded tail, so it is
unbounded.
Global leastness of $n$ gives $z\geq n$, and (8) proves (13).  Finally,
$e_x>1$ and $x\geq2n+1$ imply
\[
 \mathcal KQ=x+e_x>2(n+1),
\]
which is (14).
\end{proof}

\begin{theorem}[exact reset renewal on a $C$-free tail]
Assume that, from some record onward, no atom $C$ occurs.  A \emph{reset}
is a record immediately following an atom $B$ or $D$.  There are
infinitely many resets.  Between two consecutive resets the atom word is
uniquely of one of the forms
\[
 A^kB,
 \qquad
 A^kD,
 \qquad k\geq0.
 \tag{17}
\]

Let $x$ and $x'$ be the states at two consecutive resets, and put
\[
 Y:=x+1,
 \qquad
 Y':=x'+1,
 \qquad
 \nu:=v_2(Y),
 \qquad
 u:=\frac{Y}{2^\nu}.
 \tag{18}
\]
Then $u$ is odd, $3\nmid Y$, $3\nmid Y'$, and the following two exact
alternatives hold:
\[
\boxed{
\begin{array}{c|c|c|c}
\text{reset block}&\nu&Y'&Q'/Q\\ \hline
A^kB&k+2&\displaystyle\frac{3^\nu u+1}{2}
&\displaystyle\frac{3^\nu}{2^{\nu+1}}\\[8pt]
A^kD&k+3&\displaystyle\frac{3^{\nu+1}u+7}{8}
&\displaystyle\frac{3^{\nu+1}}{2^{\nu+3}}.
\end{array}}
\tag{19}
\]
If $i$ is the record index of the first reset in (19), then, with
$\alpha=\log_2 3$,
\[
\boxed{
 \nu\leq(\alpha-1)i+\log_2(\mathcal KQ_{j_0}).}
\tag{20}
\]
\end{theorem}

\begin{proof}
By (5), the atoms $B$ and $D$ end at decoder depth zero, so every reset
satisfies $3\nmid Y$.  An eventual tail of atoms $A$ alone would give an
infinite exact exponent-one run.  This is impossible because the exact
run from a positive odd state has length $v_2(x+1)-1$.  Hence $B$ or $D$
occurs infinitely often, proving that there are infinitely many resets and
that consecutive resets have one of the forms in (17).

The raw exponent words in (17) are
\[
 1^{k+1}2,
 \qquad
 1^{k+2}22,
 \tag{21}
\]
respectively.  The exact initial exponent-one runs therefore have lengths
$k+1$ and $k+2$.  The exponent-one tower formula gives
$\nu=k+2$ in the first case and $\nu=k+3$ in the second.

Across $A^kB$, the equations for $Y=x+1$ give
\[
 2^{k+3}Y'=3^{k+2}Y+2^{k+2}.
 \tag{22}
\]
Substitute $Y=2^{k+2}u$ and divide by $2^{k+2}$ to obtain the first state
formula in (19).  The coefficient multiplier is
\[
 \left(\frac32\right)^k\frac98
 =\frac{3^{k+2}}{2^{k+3}},
\]
which is the first multiplier in (19).

Across $A^kD$, the corresponding equation is
\[
 2^{k+6}Y'=3^{k+4}Y+7\,2^{k+3}.
 \tag{23}
\]
Substitution of $Y=2^{k+3}u$ proves the second state formula.  Its
coefficient multiplier is
\[
 \left(\frac32\right)^k\frac{81}{64}
 =\frac{3^{k+4}}{2^{k+6}},
\]
which proves the remaining entry in (19).  The endpoint is again a reset,
so $3\nmid Y'$.

Finally, $2^\nu\mid Y$ and $Y<\mathcal KQ_{j_i}$.  Consecutive record
multipliers are at most $3/2$, hence
\[
 Q_{j_i}\leq Q_{j_0}\left(\frac32\right)^i.
\]
Taking binary logarithms proves (20).
\end{proof}

\begin{corollary}[coalescing versus two-map renewal]
Every eventual four-atom interior tail has at least one of the following
two properties.
\begin{enumerate}
 \item The atom $C=1113$ occurs infinitely often.  Every such occurrence
 produces the exact half-size coalescing seed in (8), and on a globally
 least unbounded ray it can occur only after the coefficient threshold
 (14).
 \item The atom $C$ occurs only finitely often.  After its last occurrence,
 the orbit has infinitely many decoder-zero resets governed exactly by the
 two renewal laws (19), subject to the valuation budget (20).  Its atom
 word still satisfies the fixed-seed return-lattice and factor-complexity
 bounds for an interior record-atom tail.
\end{enumerate}
\end{corollary}

\begin{proof}
The two alternatives distinguish infinite from finite occurrence of
$C$.  Theorem~1 proves the first case.  After the last occurrence,
Theorem~2 supplies the reset renewal.

The fixed-seed return-lattice theorem applies to every interior
four-atom tail and gives the final assertion.
\end{proof}

\end{document}
